\documentclass[10pt]{article}
\usepackage[utf8]{inputenc}
\usepackage[french]{babel}
\usepackage[T1]{fontenc}
\usepackage{geometry}
\usepackage{array, booktabs, makecell, xcolor}
\usepackage{amsmath, amssymb}
\usepackage{graphicx}
\geometry{a4paper, margin=15mm, landscape}
\renewcommand{\arraystretch}{1.4}
\setlength{\tabcolsep}{5pt}

\title{\Huge Fiche de préparation}
\author{\Large Timothée Meneux, Benjamin Wojtecki}
\date{}

\begin{document}
\maketitle

\rowcolors{3}{gray!10}{white}

\begin{center}
\renewcommand\cellalign{tl}
\begin{tabular}{|p{3.5cm}|p{6.5cm}|p{3.5cm}|p{3.5cm}|p{2cm}|}
\hline
\multicolumn{5}{|c|}{\textbf{\Large Timothée Meneux \& Benjamin Wojtecki}} \\
\hline
\textbf{Classe :} CM2 & \multicolumn{3}{l|}{\textbf{Titre : La ville embourbée}} & \textbf{Séance n°3} \\
\hline
\multicolumn{5}{|l|}{\textbf{Compétences travaillées :} Communication, élaboration d'un vocabulaire} \\
\multicolumn{5}{|l|}{\textbf{Objectifs :} Appréhender la notion d'arbre couvrant} \\
\hline
\rowcolor{gray!30}
\textbf{Activité} & \textbf{Script} & \textbf{Notes / Objectifs} & \textbf{Matériel} & \textbf{Temps} \\
\hline
Présentation de l'activité 
& \raggedright Introduction devant la classe entière.\\ Histoire et consignes. 
& & & 5 min \\
\hline
Distribution du matériel 
& \raggedright Donner la première fiche. 
& & Les différentes cartes & 2 min \\
\hline
Activité 
& \raggedright Recherche par groupe de 2 ou 3.\\ Extension en donnant des couleurs aux maisons.
& \raggedright Trouver l'algorithme de Kruskal.\\ Comprendre la résistance aux changements.
& & 20 min \\
\hline
Remise en commun 
& \raggedright Explication de l'algorithme de Kruskal avec un schéma.
& \raggedright Expliquer l'algorithme de Kruskal.
& Dessin à la main ou animation & 10 min \\
\hline
Application d'un algorithme 
& \raggedright On leur donne un gros graphe et ils appliquent l'algorithme.
& \raggedright Augmentation de la difficulté.
& & 5 min \\
\hline
Institutionnalisation 
& \raggedright Une application pratique de l’algorithme de Kruskal consiste à relier tous les ordinateurs d’un réseau local en optimisant la longueur de fibre optique.\\
On peut aussi l’utiliser pour déterminer le chemin optimal entre plusieurs villes.
& & & 8 min \\
\hline
\multicolumn{4}{|r|}{\textbf{Total :}} & \textbf{50 min} \\
\hline
\end{tabular}
\end{center}

\end{document}
